\chapter{Краткий исторический очерк}

\section{Гуанчи}

До~европейского завоевания остров населяли \place{гуанчи} ---
берберского происхождения племена, пришедшие с~материка около
500~г. до~н.\,э. Они жили в~пещерах, занимались земледелием
(ячмень, гофио), разводили коз; управлялись двумя вождями
(\textit{гуанартеме}) --- в~Тельде и~Гальдаре. Следы их~быта
сохранились в~\sight{Куэва-Пинтада} (Гальдар) и~в~пещерах
\sight{Барранко-де-Гуаядеке}.

\marginfact{Покорение\\острова:\\1478–1483}

\section{Кастильское завоевание}

В~1478\,г. Хуан Рехон высадился у~устья Гвинигуады и~основал
первый форт --- \placex{Реаль-де-Лас-Пальмас}. После пятилетнего
сопротивления, закончившегося падением крепости Ансите и~легендой
о~самоубийственном прыжке последних воинов в~пропасть
\textit{Fortaleza de Ansite}, остров был включён в~Кастильскую
корону (1483).

\section{Сахар, кошениль, бананы}

XVI~в. --- эпоха сахарного тростника; XVII–XVIII~вв. --- виноделия
(\textit{canary sack}, известного Шекспиру); XIX~в. --- кошенили,
дававшей красную краску; с~конца XIX~в. доминируют бананы
и~трансатлантическая торговля углём.

\section{Туризм}

Массовый туризм начался в~1960-е гг.\ с~освоением южного побережья
около \place{Маспаломас} и~\place{Плайя-дель-Инглес}. Сегодня это
главная статья экономики острова, однако внутренние горные
муниципалитеты сохраняют сельский, почти средневековый уклад.
